\begin{textsample}{NEG Dim 3 – 2021\_c6 – Score 3.00 – 2021\_c6/enade\_2021\_ci\_bio\_d313503e1f.txt}  \label{ex:f3_neg_242}
CAPÍTULO 6 ANÁLISE TÉCNICA DA PROVA Este capítulo tem por objetivo apresentar o desempenho dos estudantes concluintes de Ciências Biológicas ( Licenciatura/Bacharelado ) no Enade/2021.
Para isso, foram calculadas as estatísticas básicas da prova em seu todo ( seção 6.1.1 ), bem como as estatísticas relacionadas aos Componentes de Formação Geral ( seção 6.1.2 ) e de Conhecimento Específico da Área ( seção 6.1.3 ).
Dadas as suas características, foram analisadas, em separado, as questões objetivas ( seção 6.2 ) e as questões discursivas ( seção 6.3 ).
Tomando -se como base as duas Questões Discursivas do Componente de Formação Geral, nas seções 6.3.1.6 a 6.3.1.8, são apresentados os resultados e os comentários sobre a correção das respostas em relação à Língua Portuguesa.
Para as questões objetivas, foram disponibilizados os Índices de Facilidade e de Discriminação Ponto-Bisserial, também em separado, para os Componentes de Formação Geral ( seção 6.2.1 e 6.2.2 ) e de Conhecimento Específico ( seção 6.2.3 e 6.2.4 ).
De cada componente, uma das questões foi escolhida para exemplificar a análise gráfica, relacionando -se as alternativas escolhidas pelos estudantes ( inclusive o gabarito ) com o número de acertos no componente.
Nos Anexos IA e IB, apresenta -se a íntegra da análise gráfica para todas as questões objetivas, respectivamente à Licenciatura e ao Bacharelado.
Para cada uma das questões discursivas, os conteúdos dos tipos mais comuns de respostas dos estudantes são apresentados e comparados com o padrão de resposta esperado ( ver Anexos IXA e IXB com o padrão de respostas ).
Nas tabelas, constam as seguintes estatísticas das notas¹⁹ : média do desempenho na prova, erro padrão da média, desvio padrão, nota mínima, mediana e nota máxima para cada um de seus componentes.
Tais estatísticas contemplam o total de estudantes concluintes da Área de Ciências Biológicas ( Licenciatura/Bacharelado ) inscritos e presentes à prova do Enade/2021, tendo em vista agregações, ou por Grandes Regiões e o país como um todo, segundo Habilitação e Modalidade de Ensino, ou por Categoria Administrativa, Organização Acadêmica da IES.
Em relação aos gráficos de distribuição de notas, o intervalo considerado foi de 10 unidades, aberto à esquerda e fechado à direita, com exceção do primeiro intervalo, [ 0 ; 10 ], fechado em ambos os extremos.
Para os gráficos de distribuição das notas das questões discursivas, foram consideradas mais duas categorias : questão em branco²⁰ e nota zero. 19 Uma definição dessas estatísticas pode ser encontrada no Glossário. 20 Nesse grupo estão incluídas também as respostas classificadas como nulas ou desconsideradas. 6.1 ESTATÍSTICAS BÁSICAS DA PROVA Esta seção apresenta estatísticas selecionadas e histogramas da nota geral ( 6.1.1 ) e de cada componente : Formação Geral ( 6.1.2 ) e Conhecimento Específico ( 6.1.3 ).
São também apresentadas estatísticas selecionadas de subpopulações, caracterizadas por Grande Região, Categoria Administrativa, Organização Acadêmica e Modalidade de Ensino. 6.1.1 ESTATÍSTICAS BÁSICAS GERAIS Na Tabela 6.1, são apresentadas as Estatísticas Básicas da prova, por Grande Região segundo Habilitação e Modalidade de Ensino, dos estudantes concluintes de Ciências Biológicas ( Licenciatura/Bacharelado ).
As análises serão feitas primeiramente para os cursos de Licenciatura e depois para os de Bacharelado.
Cada Habilitação é desagregada internamente em cursos Presenciais e cursos a Distância.
A Média das notas da prova, como um todo, dos alunos de Licenciatura ( nas seções seguintes serão analisados os Componentes de Formação Geral e de Conhecimento Específico ), foi 42,9, sendo que os estudantes da região Norte obtiveram a Média mais baixa ( 38,0 ), e os da região Sudeste obtiveram a Média mais alta ( 45,3 ).
As demais Médias foram : 40,7, na região Nordeste, 43,9, na região Sul e 42,2, na região Centro-Oeste.
O Desvio padrão para o Brasil, como um todo, foi 13,8, sendo o maior Desvio padrão encontrado na região Sudeste ( 14,2 ), e o menor, na região Norte ( 12,4 ), indicando uma dispersão um pouco menor das notas dessa última região.
A região que obteve a maior nota Máxima foi a Sul ( 86,4 ), ao passo que a região que atingiu a menor nota Máxima foi a Norte ( 77,9 ).
A Mediana do Brasil, como um todo, foi 42,7, sendo a maior Mediana obtida na região Sudeste ( 45,6 ), e a menor obtida na região Norte ( 37,7 ).
A nota Mínima do Brasil foi obtida nas regiões Nordeste e Sudeste, onde foi zero.
A maior nota Mínima, entre as cinco Grandes Regiões, foi 7,2, na região Norte, as demais regiões obtiveram as seguintes notas Mínima : 3,2 ( Sul ) e 6,4 ( Centro-Oeste ).
Considerando -se as notas, como um todo dos alunos de Licenciatura segundo Grande Região, observa -se que existe diferença estatisticamente significativa ao nível de 95 \% entre as Médias obtidas em todas as regiões.
A Média das notas para os alunos de cursos Presenciais de Licenciatura, foi 44,1, sendo que os alunos da região Norte obtiveram a Média mais baixa ( 39,0 ), e os da região Sudeste obtiveram a Média mais alta ( 48,5 ).
As demais Médias foram : 41,3 na região Nordeste, 47,0 na região Sul e 42,6 na região Centro-Oeste.
O Desvio padrão das Licenciaturas Presenciais foi 13,5, sendo o menor obtido na região Norte ( 12,5 ) e o maior na região Centro-Oeste ( 13,9 ).
A maior Mediana foi a da região Sudeste ( 49,2 ) e a menor foi a da região Norte ( 38,6 ).
A maior nota Máxima foi 86,4, obtida nas regiões Sudeste e Sul, registrada como a nota Máxima de todas as Licenciaturas.
As notas Mínimas foram diferentes de zero nas regiões Norte ( 7,2 ), Sul ( 6,6 ) e Centro-Oeste ( 6,4 ).
Na Modalidade de Educação Presencial de Licenciatura as diferenças entre quaisquer das Médias obtidas nas Grandes Regiões são estatisticamente significativas.
A Média das notas da prova dos alunos de cursos a Distância de Licenciatura, foi 39,0, mais baixa do que os alunos de cursos Presenciais ( 44,1 ).
O Desvio padrão foi 13,8.
Para esta combinação de Modalidade e Habilitação, a nota Máxima foi 83,7, a Mediana foi 38,0 e a nota Mínima foi 3,2.
Cumpre notar que ainda que as notas dos cursos Presenciais sejam comparáveis com as dos cursos a Distância dentro de uma mesma Habilitação, não é possível comparar as notas obtidas pelos alunos de Licenciatura e de Bacharelado, já que as provas são, em parte, diferentes.
A Média das notas da prova, como um todo dos alunos de Bacharelado ( nas seções seguintes serão analisados os Componentes de Formação Geral e de Conhecimento Específico ), foi 45,2, sendo que os estudantes da região Norte obtiveram a Média mais baixa ( 38,1 ), e os da região Sul obtiveram a Média mais alta ( 46,6 ).
As demais Médias foram : 43,9, na região Nordeste, 45,7, na região Sudeste e 43,9, nas região Centro-Oeste.
O Desvio padrão para o Brasil, como um todo, foi 12,0, sendo o maior Desvio padrão encontrado na região Centro-Oeste ( 13,0 ), e o menor, na região Norte ( 10,4 ), indicando uma dispersão um pouco menor das notas dessa última região.
A região que obteve a maior nota Máxima foi a Sudeste ( 88,9 ), ao passo que a região que atingiu a menor nota Máxima foi a Norte ( 64,0 ).
A Mediana do Brasil, como um todo, foi 45,5, sendo a maior Mediana obtida na região Sul ( 47,1 ), e a menor obtida na região Norte ( 38,2 ).
A nota Mínima do Brasil foi obtida nas regiões Nordeste e Sudeste, onde foi zero.
A maior nota Mínima, entre as cinco Grandes Regiões, foi 9,8, na região Sul, as demais regiões as notas Mínimas foram : 9,2 na região Norte e 8,6 na região Centro-Oeste.
Considerando -se as notas de todo o Brasil dos alunos de Bacharelado, segundo Grande Região, observa -se que não existe diferença estatisticamente significativa ao nível de 95 \% entre as Médias, obtidas nas regiões Nordeste e Centro-Oeste, e entre as das regiões Sudeste e Sul.
A Média das notas da prova como um todo para os alunos de cursos Presenciais de Bacharelado, foi 45,3, sendo que os alunos da região Norte obtiveram a Média mais baixa ( 38,1 ), e os da região Sul obtiveram a Média mais alta ( 47,5 ).
As demais Médias foram : 43,9 nas regiões Nordeste e Centro-Oeste, e 45,7 na região Sudeste.
A maior nota Máxima dos cursos presenciais de Bacharelado foi a da região Sudeste ( 88,9 ) e a menor foi a da região Norte ( 64,0 ).
A maior Mediana obtida na região Sul ( 48,0 ), e a menor obtida na Norte ( 38,2 ).
As notas Mínimas nos cursos Presenciais foram nas regiões Norte ( 9,2 ), Nordeste ( 0,0 ), Sudeste ( 0,0 ), Sul ( 10,2 ) e Centro-Oeste ( 8,6 ).
Considerando -se as notas segundo Grande Região, não existe diferença estatisticamente significativa ao nível de 95 \% entre as Médias obtidas nas regiões Nordeste e Centro-Oeste, e entre as das regiões Sudeste e Sul.
Somente na região Sudeste foram oferecidos cursos a Distância de Bacharelado.
A Média foi 36,5 e o Desvio padrão foi 10,4.
Para esta combinação de Modalidade e Habilitação, a nota Máxima foi 58,8, a Mediana foi 36,4 e a nota Mínima foi 9,8.
O comportamento das notas dos estudantes de todo o Brasil pode ser observado no Gráfico 6.1, que apresenta um histograma com a distribuição das mesmas, segundo a Habilitação.
Essa é uma distribuição unimodal para os alunos de cursos de Licenciatura e de Bacharelado, com ambas as modas no intervalo ( 40 ; 50 ].
Na Tabela 6.2, são apresentadas informações referentes à Média da nota final, desagregadas por Categoria Administrativa, por Organização Acadêmica, segundo a Habilitação.
A Média das notas da prova, como um todo, dos alunos de Licenciatura, foi 42,9 ( Tabela 6.1 ).
Em relação a Categoria Administrativa, os estudantes de Licenciatura das IES Públicas obtiveram Média mais alta ( 43,6 ) que a Média nacional, e os das IES Privadas obtiveram Média mais baixa ( 40,6 ).
Observa -se que existe diferença estatisticamente significativa entre as Médias das notas das IES Públicas e as das Privadas.
A diferença entre as Médias das regiões Sudeste e Norte ( 7,3 ), a maior e a menor Média, é superior à diferença entre a das IES Públicas e Privadas ( 3,0 ), não chegando a caracterizar uma maior diversidade administrativa do que regional.
O Desvio padrão para as IES Públicas ( 13,7 ) foi inferior ao do Brasil, como um todo ( 13,8 ), indicando uma dispersão um pouco menor das notas nesta Categoria Administrativa.
No tocante à Organização Acadêmica, as Universidades ( 43,9 ) obtiveram Média mais alta que a nacional ( 42,9 ).
A Média dos Centros Universitários ( 41,3 ), das Faculdades ( 37,6 ) e dos CEFET/IF ( 39,3 ) foi menor do que a nacional.
Constata -se que só não existe diferença estatisticamente significativa, ao nível de 95 \%, entre as Médias das Faculdades e CEFET/IF.
A Média das notas da prova, como um todo, dos alunos de Bacharelado, foi 45,2 ( Tabela 6.1 ).
Em relação a Categoria Administrativa, os estudantes de Bacharelado das IES Públicas obtiveram Média mais alta ( 47,9 ) que a Média nacional, e os das IES Privadas obtiveram Média mais baixa ( 41,1 ).
Observa -se que existe diferença estatisticamente significativa entre as Médias das notas das IES Públicas e as das Privadas.
A diferença entre as Médias das regiões Sul e Norte ( 8,5 ), a maior e a menor Média, é superior à diferença entre a das IES Públicas e Privadas ( 6,8 ), caracterizando -se uma maior diversidade administrativa do que regional.
O Desvio padrão para as IES Públicas ( 11,7 ) foi inferior ao do Brasil, como um todo ( 12,0 ), indicando uma dispersão um pouco menor das notas nesta Categoria Administrativa.
No tocante a Organização Acadêmica, as Universidades e os CEFET/IF obtiveram Média mais alta que a nacional ( 46,3 e 46,1, respectivamente ).
A Média dos Centros Universitários e a das Faculdades foram menores do que a nacional ( 39,8 e 38,1, respectivamente ).
Constata -se que não existe diferença estatisticamente significativa, ao nível de 95 \%, entre as Médias de Universidades e CEFET/IF, e entre as de Centros Universitários e Faculdades. 6.1.2 ESTATÍSTICAS BÁSICAS NO COMPONENTE DE FORMAÇÃO GERAL Na Tabela 6.3, são apresentadas as Estatísticas Básicas em relação ao componente da prova que avalia a Formação Geral dos estudantes concluintes.
A nota de Formação Geral é obtida levando -se em conta os acertos nas oito questões objetivas para o curso de Ciências Biológicas ( Licenciatura ), já para o curso de Ciências Biológicas ( Bacharelado ) foram levadas em conta oito questões objetivas e a nota média obtida nas Questões Discursivas 1 e 2, para ambas Habilitações, sendo a nota de cada questão discursiva obtida levando -se em conta 80 \% da nota relativa ao conteúdo e 20 \% da nota relativa ao desempenho linguístico.
Os estudantes de Licenciatura, de todo o Brasil, obtiveram desempenho médio de 37,9.
Quanto à variabilidade, o Desvio padrão das notas dos estudantes de Licenciatura do Brasil foi 15,4.
A maior Média foi obtida na região Sudeste ( 40,3 ) e a menor, na região Norte ( 33,7 ).
As demais Médias foram : 35,9, na região Nordeste, 38,9, na região Sul e 36,8, na região Centro-Oeste.
Já o maior Desvio padrão foi obtido na região Centro-Oeste ( 15,9 ), e o menor, na região Norte ( 13,9 ).
Os demais desvios foram : 14,8, na região Nordeste, 15,6, na região Sudeste e 15,2, na região Sul.
A maior nota no Componente de Formação Geral da prova do Enade/2021 foi 86,7, obtida por, pelo menos, um estudante na região Sul.
A menor nota Máxima foi obtida na região Norte ( 70,9 ).
A Mediana do Brasil, como um todo, foi 37,5, sendo a menor Mediana encontrada na região Norte ( 32,8 ), e a maior, na região Sudeste ( 40,3 ).
A nota Mínima nessa parte foi zero em todas as regiões.
Considerando -se as notas, segundo Grande Região, observa -se que não existe diferença estatisticamente significativa apenas entre as Médias das notas do Componente de Formação Geral obtidas nas regiões Nordeste e Centro-Oeste.
Já os alunos de cursos Presenciais de Licenciatura obtiveram desempenho médio de 38,8.
Quanto à variabilidade, o Desvio padrão das notas do conjunto dos estudantes das Licenciaturas Presenciais foi 15,2.
A maior Média foi obtida na região Sudeste ( 42,4 ), e a menor, na região Norte ( 34,2 ).
As demais Médias foram : 36,5 na região Nordeste, 41,4 na região Sul e 37,1 na região Centro-Oeste.
O maior Desvio padrão foi obtido na região Centro-Oeste ( 15,9 ), e o menor, nas regiões Norte ( 14,1 ).
Os demais desvios foram : 14,5 na região Sul, 14,8 na região Nordeste e 14,9 na região Sudeste.
A nota Máxima no Componente de Formação Geral da prova de alunos de cursos Presenciais de Licenciatura foi 86,5, obtida por, pelo menos, um aluno na região Sudeste.
A menor nota Máxima foi obtida na região Norte ( 70,9 ).
A Mediana do conjunto de Licenciaturas Presenciais foi 38,4, sendo a menor Mediana encontrada na região Norte ( 33,3 ), e a maior, encontrada na região Sudeste ( 42,5 ).
A nota Mínima nesta parte foi zero em todas as regiões.
Considerando -se as notas segundo Grande Região, observa -se que apenas não existe diferença estatisticamente significativa entre as Média das notas das Licenciaturas Presenciais, no Componente de Formação Geral, das regiões Nordeste e Centro-Oeste, e entre as regiões Sudeste e Sul.
A Média das notas do Componente de Formação Geral dos alunos de Licenciatura a Distância foi 35,3, mais baixa do que a dos cursos Presenciais ( 38,8 ), o mesmo ocorrendo com a Mediana ( 34,8 ).
Para esta combinação de Modalidade e Habilitação, o Desvio padrão foi 15,6, a nota Mínima foi 0,0 e a nota Máxima foi 86,7.
Os estudantes de Bacharelado, de todo o Brasil, obtiveram desempenho médio de 41,8 ( Tabela 6.3 ).
Quanto à variabilidade, o Desvio padrão das notas dos estudantes de Bacharelado do Brasil, como um todo, foi 14,9.
A maior Média foi obtida na região Sudeste ( 42,6 ), e a menor, na região Norte ( 36,5 ).
As demais Médias foram : 40,2, na região Nordeste ; 42,4, na região Sul ; e 40,6, na região Centro-Oeste.
Já o maior Desvio padrão foi obtido na região Centro-Oeste ( 15,2 ) e o menor, na região Norte ( 13,1 ).
Na região Nordeste, o Desvio padrão foi 15,1, na região Sudeste foi 14,8 e na região Sul foi 14,9.
A maior nota no Componente de Formação Geral da prova do Enade/2021 foi 90,2, obtida por, pelo menos, um estudante na região Sudeste.
A menor nota Máxima foi obtida na região Norte ( 77,7 ).
A Mediana do Brasil, como um todo, foi 42,0, sendo a menor Mediana encontrada na região Norte ( 37,5 ), e a maior, na região Sudeste ( 43,0 ).
A nota Mínima nessa parte foi diferente de zero apenas na região Centro-Oeste ( 4,8 ).
Considerando -se as notas, segundo Grande Região, observa -se que não existe diferença estatisticamente significativa entre as Médias das notas do Componente de Formação Geral obtidas nas regiões Nordeste e Centro-Oeste, como também entre as Médias da região Sul e as regiões Sudeste e Centro-Oeste.
As estatísticas das notas dos alunos de cursos Presenciais de Bacharelado das regiões Norte, Nordeste, Sudeste e Centro-Oeste são as mesmas apresentadas para o total de cursos de Bacharelados participantes do Enade em 2021.
Como a desagregação ocorre na região Sul, a maior Média foi obtida nesta região ( 43,1 ), após a desagregação, considerando -se apenas os cursos presenciais de Bacharelado.
A região que obteve a maior nota Máxima foi a Sudeste ( 90,2 ).
A Mediana do conjunto dos Bacharelados Presenciais foi 42,1, sendo a Mediana da região Sul ( 43,7 ) a maior entre as regiões.
A nota Mínima continua sendo zero nas regiões Norte, Nordeste e Sudeste, 4,8 na região Centro-Oeste, enquanto na região Sul, a nota Mínima foi 2,0.
Considerando -se as notas segundo Grande Região, observa -se que não existem diferença estatisticamente significativa ao nível de 95 \% entre as Médias das regiões Nordeste e Centro-Oeste, como também entre as Médias das regiões Sul e Sudeste.
Somente na região Sul houve participação de estudantes de cursos a Distância de Bacharelado.
Para esta combinação de Modalidade e Habilitação, a Média foi 34,2, o Desvio padrão foi 13,2, a Mediana foi 33,4, a nota Mínima foi zero e a nota Máxima foi 75,8.
O Gráfico 6.2 propicia a avaliação do desempenho dos estudantes do Componente de Formação Geral, a partir do histograma da distribuição das notas correspondentes.
A distribuição é unimodal para as notas de Licenciatura com moda em ( 30 ; 40 ], intervalo imediatamente anterior ao modal da distribuição de notas da prova, como um todo ( Gráfico 6.1 ).
Para as notas do Bacharelado, a moda ocorre no intervalo ( 40 ; 50 ], mesmo intervalo modal que o da distribuição de notas da prova, como um todo ( Gráfico 6.1 ).
Nota -se, ainda, que no Gráfico 6.2, as notas apresentam dispersão maior, quando comparadas a do Gráfico 6.1 ( distribuição das notas da prova ), confirmado pela comparação dos Desvios padrões que, em Formação Geral, aumenta de 13,8 para a nota da prova, como um todo ( Tabela 6.1 ), para 15,4 do Componente de Formação Geral ( Tabela 6.3 ), nas Licenciaturas e nos Bacharelados, aumenta de 12,0 para a nota da prova, como um todo ( Tabela 6.1 ), para 14,9 do Componente de Formação Geral ( Tabela 6.3 ).
Na Tabela 6.4, são apresentadas as informações referentes ao desempenho dos concluintes do Componente de Formação Geral, em diferentes agregações : Categoria Administrativa, Organização Acadêmica e Habilitação.
Primeiramente, com respeito aos alunos de Licenciatura, observa -se que existe diferença estatisticamente significativa entre as Médias dos tipos de Categoria Administrativa : a maior Média obtida por estudantes de Licenciatura de IES Públicas ( 38,4 ) e a menor, pelos de IES Privadas ( 36,3 ).
Nota -se que, também, não há diferença estatisticamente significativa apenas entre as Médias das Faculdades ( 33,4 ) e dos CEFET/IF ( 34,4 ).
Nas demais Organizações Acadêmicas as Médias foram : 38,8 para as Universidades e 37,4 para os Centros Universitários.
Dentre os alunos de Bacharelado, observa -se que existe diferença estatisticamente significativa entre as Médias dos tipos de Categoria Administrativa : a maior Média obtida por estudantes de Bacharelado de IES Públicas ( 44,5 ) e a menor, pelos de IES Privadas ( 37,8 ).
Nota -se que, também, só não há diferença estatisticamente significativa entre as Médias das Universidades ( 42,7 ) e dos CEFET/IF ( 42,4 ).
As demais Médias foram : 37,7 para os Centros Universitários e 34,9 para as Faculdades. 6.1.3 ESTATÍSTICAS BÁSICAS DO COMPONENTE DE CONHECIMENTO ESPECÍFICO Na Tabela 6.5, são apresentadas as Estatísticas Básicas referentes ao Componente de Conhecimento Específico da Área de Ciências Biológicas ( Licenciatura/Bacharelado ) por Grande Região, segundo Habilitação e Modalidade de Ensino.
A nota do componente de Conhecimentos Específicos leva em conta as notas da parte objetiva e a média das notas obtidas nas três questões discursivas, parte discursiva da Área.
A Média do desempenho dos estudantes de Licenciatura do Brasil, como um todo, foi 44,5.
A maior Média foi obtida na região Sudeste ( 47,0 ), e a menor, na região Norte ( 39,4 ).
As demais Médias foram : 42,3 na região Nordeste, 45,5 na região Sul e 44,0 na região Centro-Oeste.
Quanto à variabilidade das notas, o Desvio padrão do Brasil, como um todo para os alunos de Licenciatura, foi 15,2, sendo o maior Desvio padrão observado na região Sudeste ( 15,7 ) e o menor, na região Norte ( 14,2 ).
Os demais desvios foram : 14,3 na região Nordeste, 15,1 na região Sul e 15,2 na região Centro-Oeste.
A Mediana das notas dos estudantes de Licenciatura de todo o Brasil foi 44,5.
A maior Mediana ocorreu na região Sudeste ( 47,3 ), e a menor, na região Norte ( 38,4 ).
As demais Medianas foram : 42,5 na região Nordeste, 45,5 na região Sul e 43,7 na região Centro-Oeste.
A nota Máxima do Brasil, como um todo, foi 93,5, sendo obtida por pelo menos um estudante de Licenciatura na região Sul.
As demais notas Máximas foram : 85,0 na região Norte, 87,5 na região Nordeste, 91,7 na região Sudeste e 86,8 na região Centro-Oeste.
A nota Mínima foi zero nas regiões Nordeste, Sudeste e Sul.
As demais notas Mínimas foram : 4,3 na região Norte e 3,0 na região Centro-Oeste.
Observa -se que não existe diferença estatisticamente significativa apenas entre as regiões Sul e Centro-Oeste na análise das Médias das notas do Componente de Conhecimento Específico.
Já para os alunos dos cursos Presenciais de Licenciatura, a Média do desempenho de Conhecimento Específico foi 45,9.
A maior Média foi obtida na região Sudeste ( 50,5 ), e a menor, na região Norte ( 40,5 ).
As demais Médias foram : 42,9 na região Nordeste, 48,9 na região Sul e 44,4 na região Centro-Oeste.
Quanto à variabilidade das notas, o Desvio padrão do conjunto das Licenciaturas presenciais foi 15,0, sendo o maior Desvio padrão observado na região Centro-Oeste ( 15,1 ), e o menor, na região Norte ( 14,2 ).
Os demais desvios foram : 14,3 nas regiões Nordeste e Sul, e 14,8 na região Sudeste.
A Mediana das notas dos estudantes de cursos Presenciais de Licenciatura de todo o Brasil foi 46,3.
A maior Mediana ocorreu na região Sudeste ( 51,3 ), e a menor, na região Norte ( 40,2 ).
As demais medianas foram : 42,5 na região Nordeste, 49,4 na região Sul e 44,0 na região Centro-Oeste.
A nota Máxima das Licenciaturas Presenciais em Ciências Biológicas foi 93,5, sendo obtida por, pelo menos, um aluno na região Sul.
As demais notas Máximas foram : 85,0 na região Norte, 87,5 na região Nordeste, 91,7 na região Sudeste e 86,8 na região Centro-Oeste.
A nota Mínima foi diferente de zero nas regiões Norte ( 5,5 ) e Centro-Oeste ( 3,0 ).
Considerando -se as Licenciaturas oferecidas presencialmente, observa -se que existe diferença estatisticamente significativa entre as Médias das notas do Componente de Conhecimento Específico das grandes Regiões.
A Média do desempenho de Conhecimento Específico dos alunos dos cursos a Distância de Licenciatura foi 40,2, menor do que a dos cursos Presenciais.
O Desvio padrão, a Mediana, a nota Máxima e a nota Mínima foram, respectivamente, 15,2, 39,0, 88,2 e 4,3.
A Média do desempenho dos estudantes de Bacharelado do Brasil, como um todo, foi 46,3.
A maior Média foi obtida na região Sul ( 48,0 ), e a menor, na região Norte ( 38,6 ).
As demais Médias foram : 45,1, na região Nordeste ; 46,8 na região Sudeste ; e 45,1, na região Centro-Oeste.
Quanto à variabilidade das notas, o Desvio padrão do Brasil, como um todo para os alunos de Bacharelado, foi 13,0, sendo o maior Desvio padrão observado na região Centro-Oeste ( 14,1 ) e o menor, na região Norte ( 11,6 ).
Os demais desvios foram : 13,1, na região Nordeste, 12,8, na região Sudeste e 13,0, na região Sul.
A Mediana das notas dos estudantes de Bacharelado de todo o Brasil foi 46,5.
A maior Mediana ocorreu na região Sul ( 48,7 ), e a menor, na região Norte ( 39,0 ).
As demais Medianas foram : 45,3 na região Nordeste ; 47,0 na região Sudeste ; e 45,7 na região Centro-Oeste.
A nota Máxima do Brasil, como um todo, foi 92,4, sendo obtida por, pelo menos, um estudante de Bacharelado na região Sudeste.
As demais notas Máximas foram : 67,8, na região Norte ; 83,6, na região Nordeste ; 82,4, na região Sul e 78,5 na região Centro-Oeste.
A nota Mínima foi zero nas regiões Nordeste e Sudeste.
A maior nota Mínima foi obtida na região Norte ( 8,1 ).
Observa -se que não existe diferença estatisticamente significativa apenas entre as Médias, das notas do Componente de Conhecimento Específico, das regiões Nordeste e Centro-Oeste.
Já para os alunos de cursos Presenciais de Bacharelado, a Média do desempenho de Conhecimento foi ligeiramente maior, 46,4.
A maior e a menor Média são respectivamente na região Sul ( 48,9 ) e região Norte ( 38,6 ).
Na região Sul, única que teve a participação de cursos a Distância, a Média desagregada dos cursos Presenciais aumentou de 48,0 do total para 48,9.
O Desvio padrão dos cursos presenciais da região Sudeste, passou de 13,0 para 12,7, e com a Mediana, de 48,7 para 49,5.
As notas Máximas e as notas Mínimas obtidas em todas as regiões para os Bacharelados Presenciais foram as mesmas apresentadas para os Bacharelados do Brasil, como um todo.
Observa -se que, nos cursos de Bacharelados Presenciais não existe diferença estatisticamente significativa apenas entre as Médias, das notas do Componente de Conhecimento Específico, das regiões Nordeste e Centro-Oeste.
Cursos a Distância de Bacharelado em Ciências Biológicas foram ofertados somente na região Sul.
A Média foi 37,2, menor do que a dos cursos Presenciais.
O Desvio padrão, a Mediana, a nota Máxima e a nota Mínima foram, respectivamente, 11,7, 36,5, 60,7 e 8,1.
Assim como nos Gráficos 6.1 e 6.2, no Gráfico 6.3, apresentado a seguir, é possível ser feita uma avaliação do desempenho de concluintes de Ciências Biológicas, segundo a Habilitação, em relação ao Componente de Conhecimento Específico, com um histograma da distribuição das notas correspondentes.
Esta também é uma distribuição unimodal para ambas as Habilitações, e o grupo modal é o ( 40 ; 50 ], o mesmo do grupo modal da prova, como um todo ( Gráfico 6.1 ).
Quando comparado o grupo modal da Componente de Conhecimento Específico ( Gráfico 6.3 ) com grupo modal da Formação Geral ( Gráfico 6.2 ), para os estudantes de Licenciatura corresponde ao intervalo imediatamente acima e para os de Bacharelado corresponde ao mesmo intervalo.
Na Tabela 6.6, é apresentada uma comparação dos resultados em relação a Categoria Administrativa, a Organização Acadêmica e segundo a Habilitação, agora levando -se em conta o desempenho de estudantes do Componente de Conhecimento Específico da prova.
No que se refere a Organização Acadêmica, para os alunos de Licenciatura, a maior Média foi a das Universidades ( 45,6 ), vindo a seguir a dos Centros Universitários ( 42,6 ), a do CEFET/IF ( 41,0 ) e, por fim, a das Faculdades ( 38,9 ).
O maior Desvio padrão, e superior ao valor para o Brasil, como um todo, foi o dos Centros Universitários ( 15,6 ).
As Universidades obtiveram a maior nota Máxima ( 93,5 ).
Os Centros Universitários obtiveram nota Máxima 86,5, os CEFET/IF, 82,5, e as Faculdades, 81,5.
As Medianas foram : 45,8 nas Universidades, 41,7 nos Centros Universitários, 38,3 nas Faculdades e 40,5 nos CEFET/IF.
A nota Mínima foi zero para quase todas as Organizações Acadêmicas, exceto para os Centros Universitários ( 0,8 ).
Observa -se que existe diferença estatisticamente significativa, ao nível de 95 \%, no Componente de Conhecimento Específico entre as notas dos tipos de Organização Acadêmica, com exceção da diferença entre as Faculdades e os CEFET/IF.
Quanto a Categoria Administrativa, observa -se um comportamento semelhante àquele da parte de Formação Geral e à prova, como um todo, ou seja, existe diferença estatisticamente significativa entre as Médias das IES Públicas ( 45,3 ) e as das IES Privadas ( 42,0 ).
Neste caso, também, a maior Média foi obtida por estudantes de Licenciatura de IES Públicas de ensino. 6.2 ANÁLISE DAS QUESTÕES OBJETIVAS Esta seção apresenta estatísticas selecionadas e histogramas das Questões Objetivas de Formação Geral para Licenciatura ( 6.2.1 ) e para Bacharelado ( 6.2.2 ) e de Conhecimento estatisticamente significativa entre as Médias das IES Públicas ( 45,3 ) e as das IES Privadas ( 42,0 ).
Neste caso, também, a maior Média foi obtida por estudantes de Licenciatura de IES Públicas de ensino.
No que se refere a Organização Acadêmica, para os alunos de Bacharelado, a maior Média foi a das Universidades e dos CEFET/IF ( 47,4 ), vindo a seguir a dos Centros Universitários ( 40,5 ) e, por fim, a das Faculdades ( 39,1 ).
O maior Desvio padrão, abaixo ao valor para o Brasil, como um todo, foi o das Universidades ( 12,9 ).
As Universidades também obtiveram a maior nota Máxima ( 92,4 ).
Os Centros Universitários obtiveram nota Máxima 74,5 ; as Faculdades, 77,6 ; e os CEFET/IF obtiveram nota Máxima 71,9, a menor delas.
As Medianas foram : 48,0 nas Universidades, a maior, 40,5 nos Centros Universitários, 39,5 nas Faculdades, a menor delas, e 45,7, nos CEFET/IF.
A nota Mínima foi zero para as Universidades e Faculdades, 8,1 nos Centros Universitários e 26,3 nos CEFET/IF.
Observa -se que não existe diferença estatisticamente significativa, ao nível de 95 \%, no Componente de Conhecimento Específico entre as Médias das Universidades e CEFET/IF, e entre os Centros Universitários e Faculdades.
Quanto a Categoria Administrativa, observa -se um comportamento semelhante àquele da parte de Formação Geral e à prova, como um todo, ou seja, existe diferença estatisticamente significativa entre as Médias das IES Públicas ( 49,0 ) e as das IES Privadas ( 42,2 ).
Neste caso, também, a maior Média foi obtida por estudantes de IES Públicas de ensino. 6.2 ANÁLISE DAS QUESTÕES OBJETIVAS Esta seção apresenta estatísticas selecionadas e histogramas das Questões Objetivas de Formação Geral para Licenciatura ( 6.2.1 ) e para Bacharelado ( 6.2.2 ) e de Conhecimento Específico para Licenciatura ( 6.2.3 ) e para Bacharelado ( 6.2.4 ).
São, também, apresentadas e comparadas as Médias das subpopulações caracterizadas por Grande Região. 21 6.2.1 COMPONENTE DE FORMAÇÃO GERAL – CIÊNCIAS BIOLÓGICAS ( LICENCIATURA ) Na Tabela 6.7a, são apresentadas as Estatísticas Básicas relativas às oito Questões Objetivas do Componente da prova que abrange a Formação Geral dos estudantes de Licenciatura.
A Média do Brasil foi 41,1.
A menor Média foi encontrada na região Norte ( 36,8 ), e a maior, na região Sudeste ( 43,4 ).
As demais Médias foram : 39,1 na região Nordeste, 42,2 na região Sul e 40,1 na região Centro-Oeste.
O Desvio padrão do Brasil foi 19,6, sendo o maior Desvio padrão encontrado na região Centro-Oeste ( 19,9 ), e o menor Desvio padrão encontrado na região Norte ( 18,4 ).
As Medianas do Brasil, como um todo, e de todas as regiões foram iguais a 37,5.
A nota Máxima 100,0 foi alcançada em quase todas as Grandes Regiões, com exceção da região Norte ( 87,5 ).
As notas Mínimas ( 0,0 ) foram iguais para todas as regiões.
Na Tabela 6.8a, são apresentados o Índice de Facilidade e o Índice de Discriminação ( Ponto-Bisserial ) para cada uma das Questões Objetivas do Componente de Formação Geral.
Quanto ao Índice de Facilidade, foram usadas as seguintes cores para diferenciar o nível de dificuldade da questão :  Azul para as questões classificadas com índice Muito fácil ( \&gt ; =0,86 ), verde para as questões classificadas com índice Fácil ( 0,61 a 0,85 ), amarelo para as questões classificadas com Médio ( 0,41 a 0,60 ), vermelho para as questões classificadas com Difícil ( 0,16 a 0,40 ) e roxo para as questões classificadas com Muito difícil ( \&lt ; =0,15 ).
Já quanto ao Índice de Discriminação, foram usadas as seguintes cores para qualificar a questão :  As questões classificadas com índice Fraco receberam a cor vermelha ( \&lt ; =0,19 ), as classificadas com Médio receberam a cor amarela ( 0,20 a 0,29 ), as classificadas com Bom receberam a cor verde ( 0,30 a 0,39 ) e as classificadas com Muito bom ( \&gt ; =0,40 ) receberam a cor azul.
As Questões Objetivas do Componente de Formação Geral, segundo o Índice de Facilidade, foram assim avaliadas : nenhuma das oito questões foi anulada pela Comissão Assessora de Área.
Nenhuma teve o Índice de Facilidade classificado como Muito fácil, e uma questão foi tida como Fácil, por ter índice de acertos 0,63.
Quatro questões foram consideradas com índice de dificuldade Médio, situando -se no intervalo entre 0,41 e 0,59 do Índice de Facilidade, ou seja, houve entre 41,0 \% e 59,0 \% de acertos, enquanto três questões foram classificadas na categoria Difícil, apresentando acertos na faixa entre 16 \% e 40 \%.
Nenhuma questão apresentou menos de 14 \% de acertos, razão pela qual não houve questão classificada como Muito difícil.
O Índice de Facilidade variou de 0,20 a 0,63, e o de Discriminação, de 0,24 a 0,48.
Cinco questões, quanto ao Índice de Discriminação, foram tidas como Muito bom, duas questões foram tidas como Bom e uma questão foi tida como Médio. 21 As fórmulas utilizadas para o coeficiente de correlação ponto-bisserial e para o índice de facilidade estão apresentadas no capítulo 1.
Formação Geral.
Quanto ao Índice de Facilidade, foram usadas as seguintes cores para diferenciar o nível de dificuldade da questão :  Azul para as questões classificadas com índice Muito fácil ( \&gt ; =0,86 ), verde para as questões classificadas com índice Fácil ( 0,61 a 0,85 ), amarelo para as questões classificadas com Médio ( 0,41 a 0,60 ), vermelho para as questões classificadas com Difícil ( 0,16 a 0,40 ) e roxo para as questões classificadas com Muito difícil ( \&lt ; =0,15 ).
Já quanto ao Índice de Discriminação, foram usadas as seguintes cores para qualificar a questão :  As questões classificadas com índice Fraco receberam a cor vermelha ( \&lt ; =0,19 ), as classificadas com Médio receberam a cor amarela ( 0,20 a 0,29 ), as classificadas com Bom receberam a cor verde ( 0,30 a 0,39 ) e as classificadas com Muito bom ( \&gt ; =0,40 ) receberam a cor azul.
As Questões Objetivas do Componente de Formação Geral, segundo o Índice de Facilidade, foram assim avaliadas : nenhuma das oito questões foi anulada pela Comissão Assessora de Área.
Nenhuma teve o Índice de Facilidade classificado como Muito fácil, e uma questão foi tida como Fácil, por ter índice de acertos 0,63.
Quatro questões foram consideradas com índice de dificuldade Médio, situando -se no intervalo entre 0,41 e 0,59 do Índice de Facilidade, ou seja, houve entre 41,0 \% e 59,0 \% de acertos, enquanto três questões foram classificadas na categoria Difícil, apresentando acertos na faixa entre 16 \% e 40 \%.
Nenhuma questão apresentou menos de 14 \% de acertos, razão pela qual não houve questão classificada como Muito difícil.
O Índice de Facilidade variou de 0,20 a 0,63, e o de Discriminação, de 0,24 a 0,48.
Cinco questões, quanto ao Índice de Discriminação, foram tidas como Muito bom, duas questões foram tidas como Bom e uma questão foi tida como Médio.
Na Tabela 6.9a, é apresentada a distribuição das questões, levando -se em conta, simultaneamente, a classificação dos dois índices.
Como já comentado, nenhuma questão foi anulada.
As cinco questões que tiveram Índice de Discriminação Muito bom figuraram em dois níveis de dificuldade Fácil e Médio : uma classificada na categoria Fácil ( Questão 7 ) do Índice de Facilidade, e quatro, na categoria Médio ( questões 1, 3, 5 e 6 ).
Em particular, a Questão 1 foi a que apresentou o maior poder discriminatório, com índice 0,48, e foi considerada Médio em termos de Facilidade, com uma proporção de 0,43 acertos.
O máximo de acertos foi alcançado pela Questão 7, com um Índice de Facilidade de 0,63.
A Questão 8 foi considerada Difícil, com Índice de Facilidade 0,20.
O Índice de Discriminação foi Médio para a Questão 8 ( 0,24 ).
No Gráfico 6.4a, para exemplificar, analisa -se o comportamento da questão de número 3 de Formação Geral.
Trata -se de uma questão considerada com nível Médio de dificuldade, em relação à Facilidade, e uma das que obteve um dos maiores Índice de Discriminação dessa parte da prova.
Neste gráfico, cada uma das cinco curvas representa o percentual de respostas em determinada alternativa da questão, em função do número de acertos dos estudantes nessa parte da prova ( Formação Geral/Múltipla Escolha ), antes de possíveis eliminações pelo critério Ponto-Bisserial.
Como já foi mencionado, nenhuma questão objetiva do Componente de Formação Geral foi anulada pela Comissão Assessora de Área, por isso o número de acertos alcança no máximo oito para esse componente.
A curva em verde corresponde à alternativa B, a alternativa correta para essa questão.
Observa -se que entre os estudantes com menor número de acertos, nessa parte do exame, a situação mais frequente foi a escolha da alternativa E ( em vermelho ), incorreta.
Por exemplo, entre os estudantes que acertaram duas respostas, 45,1 \% escolheram a alternativa E ( em vermelho ), 28,6 \% escolheram a alternativa B ( em verde, gabarito ), 12,3 \% escolheram a alternativa D ( em laranja ), 8,6 \% escolheram a alternativa A ( em azul ) e 5,0 \% escolheram a alternativa C ( em preto ).
Dentre os que acertaram duas respostas entre as questões de múltipla escolha de Formação Geral, 0,3 \% deixou a questão em branco e 0,1 \% marcou mais de uma alternativa, invalidando a questão.
À medida que o número de acertos aumenta, indicando desempenho melhor nessa parte da prova, aumenta, concomitantemente, a proporção de estudantes que selecionaram a alternativa correta B, atingindo 100 \% para os estudantes com oito acertos.
Essa análise permite verificar como a questão discriminou os grupos de desempenho, justificando -se o alto índice ( 0,47 ) obtido na questão.
Cumpre notar que não é possível inferir deste gráfico nem o Índice de Facilidade, que seria uma média da proporção ponderada pela quantidade de estudantes com cada uma das notas, nem o Índice de Discriminação Ponto-Bisserial, por razão equivalente.
No caso extremo, no qual a grande concentração dos acertos dos estudantes fosse abaixo de quatro, o Índice de Facilidade seria obrigatoriamente abaixo de 65,0 \% ( neste exemplo ).
Caso a concentração fosse em seis acertos ou mais, o índice seria obrigatoriamente acima de 90,0 \%.
Os gráficos relativos às demais questões de Formação Geral constam do Anexo IA. 6.2.2 COMPONENTE DE FORMAÇÃO GERAL – CIÊNCIAS BIOLÓGICAS ( BACHARELADO ) Como uma questão foi anulada pela Comissão Assessora de Área, a Questão 1, na Tabela 6.7b, são apresentadas as Estatísticas Básicas relativas às sete Questões Objetivas válidas do Componente da prova que abrange a Formação Geral dos estudantes de Bacharelado.
A Média do Brasil foi 47,7.
A menor Média foi encontrada na região Norte ( 39,0 ), e a maior, na região Sul ( 48,7 ).
As demais Médias foram : 46,3, na região Nordeste, 48,4, na região Sudeste e 47,1, na região Centro-Oeste.
O Desvio padrão do Brasil foi 19,1, sendo o maior Desvio padrão encontrado na região Nordeste ( 19,7 ), e o menor, na região Norte ( 17,5 ).
Os demais desvios foram : 18,9, na região Sudeste, 19,0, na região Sul, e 19,5, na região Centro-Oeste.
As Medianas do Brasil, como um todo, e de quase todas as regiões foram iguais a 50,0, a exceção é a região Norte ( 37,5 ).
A nota Máxima 100,0 foi alcançada em quase todas as Grandes Regiões, as exceções são as regiões Norte e Centro-Oeste ( 87,5 ).
As notas Mínimas ( 0,0 ) foram iguais para todas as regiões.
Na Tabela 6.8b, são apresentados o Índice de Facilidade e o Índice de Discriminação ( Ponto-Bisserial ) para cada uma das questões objetivas do Componente de Formação Geral.
A convenção das cores é a mesma da Tabela 6.8a.
As Questões Objetivas do Componente de Formação Geral, segundo o Índice de Facilidade, foram assim avaliadas : das oito questões válidas, nenhuma teve o Índice de Facilidade classificado como Muito fácil, e duas questões foram tidas como Fácil, por ter índice de acertos entre 0,61 e 0,85.
Três questões foram consideradas com índice de dificuldade Médio, situando -se no intervalo entre 0,41 e 0,59 do Índice de Facilidade, ou seja, houve entre 41,0 \% e 59,0 \% de acertos, enquanto outras três questões foram classificadas na categoria Difícil, apresentando acertos na faixa entre 16 \% e 40 \%.
Nenhuma questão apresentou menos de 14 \% de acertos, razão pela qual não houve questão classificada como Muito difícil.
O Índice de Facilidade variou de 0,22 a 0,73 e o de Discriminação, de 0,28 a 0,46.
Seis das oito questões válidas, quanto ao Índice de Discriminação, foram tidas como Muito bom, uma questão foi tida como Bom e uma questão foi tida como Médio.
Na Tabela 6.9b, é apresentada a distribuição das questões, levando -se em conta, simultaneamente, a classificação dos dois índices.
As seis questões que tiveram Índice de Discriminação Muito bom figuraram entre três níveis de dificuldade Fácil, Médio e Difícil : duas classificadas na categoria Fácil ( questões 6 e 7 ) do Índice de Facilidade, três na categoria Médio ( questões 1, 3 e 5 ) e uma, na categoria Difícil ( Questão 4 ).
Em particular, as Questões 1 e 3 foram as que apresentaram o maior poder discriminatório, ambas com índice 0,46, e foram consideradas Médio em termos de facilidade, com uma proporção de 0,55 e 0,56 acertos, respectivamente.
O máximo de acertos foi alcançado pela Questão 7, com um Índice de Facilidade de 0,73.
A Questão 8 foi considerada Difícil, com Índice de Facilidade 0,22.
Seu Índice de Discriminação foi o mais baixo dentre as oito questões válidas : Médio ( 0,28 ).
No Gráfico 6.4b, para exemplificar, analisa -se o comportamento da Questão 3 de Formação Geral.
Trata -se de uma questão considerada Médio, em relação à Facilidade e a que obteve o maior Índice de Discriminação dessa parte da prova.
Neste gráfico, cada uma das cinco curvas representa o percentual de respostas em determinada alternativa da questão, em função do número de acertos dos estudantes nessa parte da prova ( Formação Geral/Múltipla Escolha ), antes de possíveis eliminações pelo Na Tabela 6.8b, são apresentados o Índice de Facilidade e o Índice de Discriminação ( Ponto-Bisserial ) para cada uma das questões objetivas do Componente de Formação Geral.
A convenção das cores é a mesma da Tabela 6.8a.
Na Tabela 6.9b, é apresentada a distribuição das questões, levando -se em conta, simultaneamente, a classificação dos dois índices.
As seis questões que tiveram Índice de Discriminação Muito bom figuraram entre três níveis de dificuldade Fácil, Médio e Difícil : duas classificadas na categoria Fácil ( questões 6 e 7 ) do Índice de Facilidade, três na categoria Médio ( questões 1, 3 e 5 ) e uma, na categoria Difícil ( Questão 4 ).
Em particular, as Questões 1 e 3 foram as que apresentaram o maior poder discriminatório, ambas com índice 0,46, e foram consideradas Médio em termos de facilidade, com uma proporção de 0,55 e 0,56 acertos, respectivamente.
O máximo de acertos foi alcançado pela Questão 7, com um Índice de Facilidade de 0,73.
A Questão 8 foi considerada Difícil, com Índice de Facilidade 0,22.
Seu Índice de Discriminação foi o mais baixo dentre as oito questões válidas : Médio ( 0,28 ).
No Gráfico 6.4b, para exemplificar, analisa -se o comportamento da Questão 3 de Formação Geral.
Trata -se de uma questão considerada Médio, em relação à Facilidade e a que obteve o maior Índice de Discriminação dessa parte da prova.
Neste gráfico, cada uma das cinco curvas representa o percentual de respostas em determinada alternativa da questão, em função do número de acertos dos estudantes nessa parte da prova ( Formação Geral/Múltipla Escolha ), antes de possíveis eliminações pelo critério Ponto-Bisserial. critério Ponto-Bisserial.
A curva em verde corresponde à alternativa B, a alternativa correta para essa questão.
Observa -se que, entre os estudantes com menor número de acertos, nessa parte do exame, a situação mais frequente foi a escolha da alternativa E ( em vermelho ), incorreta.
Por exemplo, entre os estudantes que acertaram duas questões, 47,9 \% escolheram a alternativa E ( em vermelho ), 27,2 \% escolheram a alternativa B ( em verde, gabarito ), 10,4 \% escolheram a alternativa A ( em azul ), 10,2 \% escolheram a alternativa D ( em laranja ) e 4,1 \% escolheram a alternativa C ( em preto ).
Dentre os que acertaram duas respostas entre as questões de múltipla escolha de Formação Geral, 0,1 \% deixou a questão em branco e 0,0 \% marcou mais de uma alternativa, invalidando a questão.
À medida que o número de acertos aumenta, indicando desempenho melhor nessa parte da prova, aumenta, concomitantemente, a proporção de estudantes que selecionaram a alternativa correta B, atingindo 100 \% para os estudantes com oito acertos.
Essa análise permite verificar como a questão discriminou os grupos de desempenho, justificando -se o alto índice ( 0,46 ) obtido na questão.
Cumpre notar que não é possível inferir deste gráfico nem o Índice de Facilidade, que seria uma média da proporção ponderada pela quantidade de estudantes com cada uma das notas, nem o Índice de Discriminação Ponto-Bisserial, por razão equivalente.
No caso extremo, no qual a grande concentração dos acertos dos estudantes fosse abaixo de quatro, o Índice de Facilidade seria obrigatoriamente abaixo de 65,0 \% ( neste exemplo ).
Caso a concentração fosse em seis acertos ou mais, o índice seria obrigatoriamente acima de 85,0 \%.
Os gráficos relativos às demais questões de Formação Geral constam do Anexo IB. 6.2.3 COMPONENTE DE CONHECIMENTO ESPECÍFICO – CIÊNCIAS BIOLÓGICAS ( LICENCIATURA ) Na Tabela 6.10a, são apresentadas as Estatísticas Básicas em relação às Questões Objetivas do Componente de Conhecimento Específico da prova de Ciências Biológicas por Grande Região.
Nessa parte da prova, como será discutido ainda nesta seção, nenhuma das Questões Objetivas do Componente de Conhecimentos Específicos foi anulada pela CAA e sete questões não foram usadas no cômputo das notas por terem sido descartadas pelo critério do Ponto-Bisserial.
Assim, as notas foram calculadas com base em 20 das 27 Questões Objetivas válidas de Conhecimentos Específicos.
A Média do Brasil deste componente foi 48,3.
A menor Média foi observada na região Norte ( 42,9 ), e a maior, na região Sudeste ( 50,7 ).
O Desvio padrão de todo o Brasil foi 16,4, sendo o menor Desvio padrão encontrado na região Norte ( 15,5 ), e o maior, na região Sudeste ( 16,8 ).
A Mediana de todo o Brasil foi 50,0, o mesmo valor da Mediana encontrada nas regiões Sudeste, Sul e Centro-Oeste.
Nas regiões Nordeste ( 45,0 ) e Norte ( 40,0 ), os valores foram menores.
Na Tabela 6.11a, são apresentados os Índices de Facilidade e de Discriminação ( Ponto-Bisserial ) das Questões Objetivas do Componente de Conhecimento Específico, para os estudantes de Ciências Biológicas de Licenciatura.
Para facilitar a diferenciação das questões, serão usadas as mesmas cores da Tabela 6.8a e 6.8b para as diferentes classificações dos Índices de Facilidade e de Discriminação.
Observa -se que nenhuma das questões foi anulada pela CAA, como já mencionado.
Sendo assim, as análises a seguir dizem respeito às 27 questões válidas.
A partir do Índice de Facilidade obtido, pode -se concluir que pouco mais da metade das questões objetivas da prova foram consideradas Difícil ou Muito difícil : das 27 questões válidas, 12 foram classificadas como Difícil e duas como Muito difícil.
Quatro questões foram classificadas como Fácil, e outras nove consideradas como Médio.
Nenhuma questão foi classificada como Muito fácil.
Já quanto aos Índices de Discriminação das Questões Objetivas do Componente de Conhecimento Específico da prova, obtém -se como resultado a seguinte classificação : quatro das 27 questões válidas foram consideradas boas, enquanto oito delas tiveram Índice de A nota Máxima da prova foi 100,00, obtida nas Questões Objetivas do Componente de Conhecimento Específico, por, pelo menos, um estudante da região Sul.
Em três regiões o valor foi 95,0, sendo elas : Nordeste, Sudeste e Centro-Oeste.
E a menor nota Máxima ocorre na região Norte ( 90,0 ).
A nota Mínima foi zero, obtida por pelo menos um estudante de quase todas regiões, com exceção da região Norte ( 5,0 ).
Na Tabela 6.11a, são apresentados os Índices de Facilidade e de Discriminação ( Ponto-Bisserial ) das Questões Objetivas do Componente de Conhecimento Específico, para os estudantes de Ciências Biológicas de Licenciatura.
Para facilitar a diferenciação das questões, serão usadas as mesmas cores da Tabela 6.8a e 6.8b para as diferentes classificações dos Índices de Facilidade e de Discriminação.
Observa -se que nenhuma das questões foi anulada pela CAA, como já mencionado.
Sendo assim, as análises a seguir dizem respeito às 27 questões válidas.
A partir do Índice de Facilidade obtido, pode -se concluir que pouco mais da metade das questões objetivas da prova foram consideradas Difícil ou Muito difícil : das 27 questões válidas, 12 foram classificadas como Difícil e duas como Muito difícil.
Quatro questões foram classificadas como Fácil, e outras nove consideradas como Médio.
Nenhuma questão foi classificada como Muito fácil.
Já quanto aos Índices de Discriminação das Questões Objetivas do Componente de Conhecimento Específico da prova, obtém -se como resultado a seguinte classificação : quatro das 27 questões válidas foram consideradas boas, enquanto oito delas tiveram Índice de Discriminação Muito bom.
Assim, para doze em 27 questões, os Índices de Discriminação foram Bom ou Muito bom.
Dentre as demais, oito delas foram classificadas como Médio, e outras sete, como Fraco, sendo 15, por conseguinte, a quantidade de questões nos dois patamares mais baixos de Discriminação.
Constata -se, assim, que a prova – no que se refere ao Componente de Conhecimento Específico – possuía média capacidade de discriminar entre aqueles que dominam ou não o conteúdo.
O Índice de Facilidade variou de 0,07 a 0,75 e o de Discriminação, de -0,01 a 0,51.
Discriminação Muito bom.
Assim, para doze em 27 questões, os Índices de Discriminação foram Bom ou Muito bom.
Dentre as demais, oito delas foram classificadas como Médio, e outras sete, como Fraco, sendo 15, por conseguinte, a quantidade de questões nos dois patamares mais baixos de Discriminação.
Constata -se, assim, que a prova – no que se refere ao Componente de Conhecimento Específico – possuía média capacidade de discriminar entre aqueles que dominam ou não o conteúdo.
O Índice de Facilidade variou de 0,07 a 0,75 e o de Discriminação, de -0,01 a 0,51.
Na Tabela 6.12a, é apresentada a distribuição das questões levando -se em conta, simultaneamente, a classificação dos dois índices.
Dentre as questões que alcançaram os maiores Índices de Discriminação, uma delas foi classificada com Índice de Discriminação Muito bom, a de número 17.
Seu índice foi 0,51, e, quanto ao Índice de Facilidade, foi A Questão de número 26 foi a mais difícil entre as 27 questões específicas válidas, com baixo Índice de Facilidade, apenas 7 \% de acertos.
Essa questão apresentou poder discriminatório muito baixo, -0,01, o que comprova ter sido a mais difícil para os estudantes.
Destaca -se, também, a Questão 21, com Índice de Facilidade 0,23, o que, em termos percentuais, corresponde a 23 \% de estudantes que responderam acertadamente, sendo, também, 0,03 seu Índice de Discriminação.
Tais questões foram, portanto, pelo critério Ponto-Bisserial, consideradas inadequadas.
Por isso, as Questões 21 e 26, além das Questões 13, 19, 22, 23 e 24, foram eliminadas do cômputo da nota final.
A título de exemplo das análises do comportamento das questões objetivas, no Gráfico 6.5a, analisa -se a Questão 17 do Componente de Conhecimento Específico.
Essa questão foi considerada, pelas respostas dos estudantes avaliados na prova, como uma questão com nível de Facilidade Médio, 0,59, ou seja, 59 \% dos estudantes assinalaram, acertadamente, a opção B correspondente ao gabarito.
Como já comentado, seu Índice de Discriminação foi igual a 0,51, classificado como Muito bom, o maior valor de Discriminação.
Neste gráfico, cada uma das cinco curvas representa o percentual de respostas em determinada alternativa da Questão 17, em função do número de acertos dos estudantes nessa parte da prova, antes de possíveis eliminações de questões pelo critério Ponto-Bisserial.
A alternativa correta B, representada no gráfico pela curva em verde, foi escolhida em maiores proporções pelos estudantes com desempenho melhor nessa parte da prova.
Já as alternativas incorretas, também denominadas distratores, foram selecionadas, principalmente, por aqueles com notas mais baixas.
No caso de Ciências Biológicas ( Licenciatura ), como se observa no eixo horizontal do Gráfico 6.5a, nenhum estudante acertou mais do que 24 questões dentre as 27 questões objetivas válidas, considerando -se o gabarito original.
A eliminação das sete questões com Índice de Discriminação Fraco do cômputo da nota é que possibilitou a existência de nota Máxima igual a 100,0 por pelo menos um estudante da região Sul ( visto na Tabela 6.10a ).
Observa -se que a soma não é 100 \%, por causa das questões não respondidas ou com mais de uma opção marcada.
Aqueles com nota zero, na sua quase totalidade, marcaram a opção A ( linha azul ).
A proporção de estudantes que selecionaram a resposta correta A aumenta gradativamente, chegando a atingir 100 \% para 22 acertos ou mais, enquanto a proporção dos que escolheram alternativas incorretas decai, a partir do primeiro acerto, em função do número de acertos nessa parte da prova.
Os gráficos relativos às demais questões do Conhecimento Específico constam do Anexo IA. 6.2.4 COMPONENTE DE CONHECIMENTO ESPECÍFICO – CIÊNCIAS BIOLÓGICAS ( BACHARELADO ) Na Tabela 6.10b, são apresentadas as Estatísticas Básicas em relação às Questões Objetivas do Componente de Conhecimento Específico da prova de Ciências Biológicas por Grande Região.
Nessa parte da prova, nenhuma questão foi anulada pela Comissão Assessora da Área.
Além disso, como será discutido ainda nesta seção, seis questões objetivas do Componente de Conhecimentos Específicos não foram usadas no cômputo das notas por terem sido descartadas pelo critério do Ponto-Bisserial.
Assim, as notas foram calculadas com base em 21 das 27 Questões Objetivas válidas de Conhecimentos Específicos.
A Média do Brasil deste componente foi 48,9.
A menor Média foi observada na região Norte ( 41,6 ), e a maior, na região Sul ( 50,5 ).
O Desvio padrão de todo o Brasil foi 13,7, sendo o menor Desvio padrão encontrado na região Norte ( 12,6 ), e o maior, na região Centro-Oeste ( 14,9 ).
A Mediana de todo o Brasil foi 47,6, o mesmo valor da Mediana encontrada nas regiões Nordeste, Sudeste e Centro-Oeste.
Na região Norte ( 42,9 ) a Mediana foi a menor, e na região Sul ( 52,4 ) foi a maior.
A nota Máxima da prova foi 95,2, obtida nas Questões Objetivas do Componente de Conhecimento Específico, por, pelo menos, um estudante da região Sudeste.
A nota Mínima foi zero nas regiões Nordeste, Sudeste e Sul.
Na Tabela 6.11b, são apresentados os Índices de Facilidade e de Discriminação ( Ponto-Bisserial ) das Questões Objetivas do Componente de Conhecimento Específico, para os estudantes de Ciências Biológicas de Bacharelado.
Para facilitar a diferenciação das questões, serão usadas as mesmas cores da Tabela 6.8a, 6.8b e 6.11a para as diferentes classificações dos Índices de Facilidade e de Discriminação.
Observa -se que nenhuma questão foi anulada pela Comissão Assessora da Área ( CAA ) e que as análises dizem respeito às 27 questões válidas.
A partir do Índice de Facilidade obtido, pode -se concluir que mais da metade das questões objetivas da prova foram considerados Difícil ( 12 ) ou Muito difícil ( três ), 15 das 27 notas por terem sido descartadas pelo critério do Ponto-Bisserial.
Assim, as notas foram calculadas com base em 21 das 27 Questões Objetivas válidas de Conhecimentos Específicos.
A Média do Brasil deste componente foi 48,9.
A menor Média foi observada na região Norte ( 41,6 ), e a maior, na região Sul ( 50,5 ).
O Desvio padrão de todo o Brasil foi 13,7, sendo o menor Desvio padrão encontrado na região Norte ( 12,6 ), e o maior, na região Centro-Oeste ( 14,9 ).
A Mediana de todo o Brasil foi 47,6, o mesmo valor da Mediana encontrada nas regiões Nordeste, Sudeste e Centro-Oeste.
Na região Norte ( 42,9 ) a Mediana foi a menor, e na região Sul ( 52,4 ) foi a maior.
A nota Máxima da prova foi 95,2, obtida nas Questões Objetivas do Componente de Conhecimento Específico, por, pelo menos, um estudante da região Sudeste.
A nota Mínima foi zero nas regiões Nordeste, Sudeste e Sul.
Na Tabela 6.11b, são apresentados os Índices de Facilidade e de Discriminação ( Ponto-Bisserial ) das Questões Objetivas do Componente de Conhecimento Específico, para os estudantes de Ciências Biológicas de Bacharelado.
Para facilitar a diferenciação das questões, serão usadas as mesmas cores da Tabela 6.8a, 6.8b e 6.11a para as diferentes classificações dos Índices de Facilidade e de Discriminação.
Observa -se que nenhuma questão foi anulada pela Comissão Assessora da Área ( CAA ) e que as análises dizem respeito às 27 questões válidas.
A partir do Índice de Facilidade obtido, pode -se concluir que mais da metade das questões objetivas da prova foram considerados Difícil ( 12 ) ou Muito difícil ( três ), 15 das 27 questões válidas.
Uma questão foi classificada como Muito fácil, seis questões foram classificadas como Fácil, e outras cinco consideradas como Médio.
Já quanto aos Índices de Discriminação das questões objetivas do Componente de Conhecimento Específico da prova, obtém -se como resultado a seguinte classificação : sete das 27 questões foram consideradas como Bom, enquanto quatro delas tiveram Índice de Discriminação Muito bom.
Assim, para 11 em 27 questões, os Índices de Discriminação foram Bom ou Muito bom.
Dentre as demais, dez delas foram classificadas como Médio, e seis, como Fraco, sendo 16, por conseguinte, a quantidade de questões nos dois patamares mais baixos de discriminação.
Constata -se, assim, que a prova – no que se refere ao Componente de Conhecimento Específico – possuía média capacidade de discriminar entre aqueles que dominam ou não o conteúdo.
O Índice de Facilidade variou de 0,08 a 0,88 e o de Discriminação, de -0,03 a 0,49.
Na Tabela 6.12b, é apresentada a distribuição das questões levando -se em conta, simultaneamente, a classificação dos dois índices.
Dentre as questões que alcançaram os maiores Índices de Discriminação, uma delas obteve o maior deles, a de número 17.
Seu índice foi 0,49, e, quanto ao Índice de Facilidade, foi classificada como Fácil, com 75 \% dos estudantes marcando a opção correta.
A Questão de número 26 foi a mais difícil entre as 27 questões específicas válidas, com baixo Índice de Facilidade, apenas 8 \% de acertos.
Essa questão apresentou poder discriminatório muito baixo, 0,08, o que comprova ter sido a mais difícil para os estudantes.
Destaca -se, também, a Questão 30, com Índice de Facilidade 0,15, o que, em termos percentuais, corresponde a 15 \% de estudantes que responderam acertadamente.
Já -0,03 foi o seu Índice de Discriminação.
Tais questões foram, portanto, pelo critério Ponto-Bisserial, consideradas inadequadas.
Por isso, as Questões 26 e 30, além das Questões 13, 19, 21 e 33, foram eliminadas do cômputo da nota final.
A título de exemplo das análises do comportamento das questões objetivas, no Gráfico 6.5b, analisa -se a Questão 17 do Componente de Conhecimento Específico.
Essa questão foi considerada, pelas respostas dos estudantes avaliados na prova, como uma questão com nível de facilidade Fácil, 0,75, ou seja, 75 \% dos estudantes assinalaram, acertadamente, a opção B ( em verde ), correspondente ao gabarito.
Como já comentado, seu Índice de Discriminação foi igual a 0,49, classificado como Muito bom, o maior valor de discriminação.
Neste gráfico, cada uma das cinco curvas representa o percentual de respostas em determinada alternativa da Questão 17, em função do número de acertos dos estudantes nessa parte da prova, antes de possíveis eliminações de questões pelo critério Ponto-Bisserial.
A alternativa correta B, representada no gráfico pela curva em verde, foi escolhida em maiores proporções pelos estudantes com desempenho melhor nessa parte da prova.
Já as alternativas incorretas, também denominadas distratores, foram selecionadas, principalmente, por aqueles com notas mais baixas.
No caso de Ciências Biológicas ( Bacharelado ), como se observa no eixo horizontal do Gráfico 6.5b, nenhum estudante acertou mais do que 22 questões dentre as 27 questões objetivas válidas, considerando -se o gabarito original.
Mesmo com a eliminação das seis questões com Índice de Discriminação Fraco do cômputo da nota, nenhum participante atingiu a nota Máxima, igual a 100,0, nessa parte da prova ( visto na Tabela 6.10b ).
Observa -se que a soma não é 100 \%, por causa das questões não respondidas ou com mais de uma opção marcada.
Aqueles com nota zero, na sua quase totalidade, marcaram a opção D ( em laranja ).
A proporção de estudantes que selecionaram a resposta correta B aumenta gradativamente, chegando a atingir 100 \% para 19 acertos ou mais, enquanto a

% matched lemmas: 
\end{textsample}
